\section{Background}
\label{sec:background}

This section fixes notation and states the three ingredients the rest of the
paper combines: tensor decompositions, the Mixture-of-Experts layer viewed as
a tensor, and the known solution of modular addition.

\subsection{Tensors and the mode-$n$ product}

A tensor of order $d$ is an element $\mathcal{W} \in \R^{I_1 \times \cdots
\times I_d}$; its $n$-th index is its $n$-th \emph{mode}. Two operations are
used throughout. The mode-$n$ unfolding $\mathcal{W}_{(n)} \in \R^{I_n \times
\prod_{m \neq n} I_m}$ rearranges $\mathcal{W}$ into a matrix whose rows are
indexed by mode $n$. The mode-$n$ product with a matrix $A \in \R^{J \times
I_n}$, written $\mathcal{W} \times_n A$, contracts the $n$-th mode:
\begin{equation}
  (\mathcal{W} \times_n A)_{i_1 \ldots j \ldots i_d}
  = \sum_{i_n = 1}^{I_n} \mathcal{W}_{i_1 \ldots i_n \ldots i_d} A_{j i_n}.
\end{equation}
Contracting a mode with a vector removes that mode. This is the operation that
makes expert layers computable without materialising them: routing weights
enter as vectors contracted against the expert mode.

We use the following standard decompositions; see \citet{kolda2009tensor} for
a full treatment.

\paragraph{CP.} The canonical polyadic decomposition writes a tensor as a sum
of $R$ rank-one terms,
\begin{equation}
  \mathcal{W} \approx \sum_{r=1}^{R} a^{(1)}_r \otimes a^{(2)}_r \otimes
  \cdots \otimes a^{(d)}_r,
  \qquad a^{(n)}_r \in \R^{I_n},
\end{equation}
costing $R \sum_n I_n$ parameters. It is maximally economical but rigid: one
rank index is shared by every mode, so the modes cannot interact except
through $r$.

\paragraph{Tucker.} The Tucker decomposition retains a dense core tensor and
compresses each mode separately,
\begin{equation}
  \mathcal{W} \approx \mathcal{Z} \times_1 A^{(1)} \times_2 A^{(2)} \cdots
  \times_d A^{(d)},
  \qquad \mathcal{Z} \in \R^{R_1 \times \cdots \times R_d},
  \quad A^{(n)} \in \R^{I_n \times R_n},
\end{equation}
costing $\prod_n R_n + \sum_n I_n R_n$ parameters. The core grows
exponentially in the order, which is the price of allowing every pair of modes
to interact. CP is the special case in which the core is diagonal.

\paragraph{Tensor-Train and Tensor-Ring.} The Tensor-Train decomposition
\citep{oseledets2011tt} arranges the modes in a chain of third-order cores,
\begin{equation}
  \mathcal{W}_{i_1 \ldots i_d} \approx
  \sum_{r_1 \ldots r_{d-1}}
  G^{(1)}_{i_1 r_1} G^{(2)}_{r_1 i_2 r_2} \cdots G^{(d)}_{r_{d-1} i_d},
\end{equation}
costing $O(d I R^2)$ parameters, linear in the order rather than exponential.
Tensor-Ring generalises it by closing the chain, so that the first and last
cores share a rank index $R_1 > 1$ and the contraction becomes a trace.

\paragraph{Hierarchical Tucker.} Instead of a chain, the modes may be arranged
on a binary tree, with transfer tensors at internal nodes
\citep{grasedyck2010hierarchical}. This retains the logarithmic depth of a
hierarchy while keeping parameter count linear in the order. We do not train a
hierarchical Tucker layer here; it is the natural extension of this study and
is discussed in Section~\ref{sec:future}.

\subsection{Mixture-of-Experts layers as tensors}

A Mixture-of-Experts layer holds $N$ experts $E_1, \ldots, E_N$ and a router
that produces weights $\gvec \in \R^N$, most of them zero, so that the layer
computes $\sum_{i} \gvec_i E_i(x)$. When each expert is linear, its parameters
form a matrix $W_i \in \R^{d_{\mathrm{out}} \times d_{\mathrm{in}}}$, and
stacking them gives a third-order tensor
\begin{equation}
  \mathcal{W} \in \R^{d_{\mathrm{out}} \times d_{\mathrm{in}} \times N},
  \qquad
  y = \mathcal{W} \times_2 x \times_3 \gvec .
  \label{eq:moe-tensor}
\end{equation}
Equation~\eqref{eq:moe-tensor} is the unifying view of this paper: an
architecture is a choice of how to represent $\mathcal{W}$, and every layer we
compare is one such choice.

\paragraph{Product-key routing.} With $N$ experts a dense router costs $O(N)$.
Product keys \citep{lample2019productkeys} avoid this by splitting the input
into halves $x^{(1)}, x^{(2)}$, matching each against $n = \sqrt{N}$ sub-keys
per head $h$, and indexing experts by the pair $(i,j)$:
\begin{equation}
  z^{(1)}_{hi} = \langle w^{(1)}_{hi}, x^{(1)} \rangle, \quad
  z^{(2)}_{hj} = \langle w^{(2)}_{hj}, x^{(2)} \rangle, \quad
  \gvec^{(1)}_h = \softmax\!\big(\TopK(z^{(1)}_h)\big), \quad
  \gvec_{hij} = \gvec^{(1)}_{hi}\,\gvec^{(2)}_{hj}.
  \label{eq:product-key}
\end{equation}
The weight on expert $(i,j)$ is thus a product of two independent choices,
which is itself a rank-one constraint on the routing distribution.

\paragraph{The architectures compared.} Splitting the expert index into two
modes turns $\mathcal{W}$ into a fourth-order tensor in $\R^{d_{\mathrm{out}}
\times d_{\mathrm{in}} \times n \times n}$, and the layers we compare are
different representations of it.

\emph{Unfactorised.} Every expert is stored separately, at cost $O(N \cdot
d_{\mathrm{out}} d_{\mathrm{in}})$. No structure is shared.

\emph{MONET-HD} \citep{park2025monet}. The expert is split horizontally into a
bottom factor indexed by $i$ and a top factor indexed by $j$,
\begin{equation}
  E_{ij}(x) = V_j\, \sigma\!\big(U_i x + b^{(1)}_i\big) + b^{(2)}_j,
  \qquad U_i \in \R^{m \times d_{\mathrm{in}}},\ V_j \in \R^{d_{\mathrm{out}}
  \times m},
  \label{eq:hd}
\end{equation}
so $2n$ factors generate $n^2$ experts and the parameter count grows as
$\sqrt{N}$. Unlike the other layers here this one is not multilinear: the
nonlinearity $\sigma$ sits between the two factors.

\emph{MONET-VD.} The expert is instead split vertically, with the first half
of the output produced from factor $i$ and the second half from factor $j$,
each acting on a concatenation of both halves of the hidden state. The
parameter count matches HD.

\emph{CP, Tucker, TT.} Following \citet{oldfield2024mumoe}, the fourth-order
tensor is given a CP, Tucker or Tensor-Train representation directly, with the
two expert modes contracted against $\gvec^{(1)}_h$ and $\gvec^{(2)}_h$. These
layers are multilinear: the output is a product of a factor depending on $x$
and factors depending on each routing vector, with no elementwise
nonlinearity between them.

Table~\ref{tab:params} gives parameter counts for $P = 113$, $m = 8$ and four
router heads. The contrast that matters is the growth rate: as $N$ increases
by a factor of $256$, the unfactorised layer grows by $238\times$, the MONET
layers by exactly $16 = \sqrt{256}$, and CP and Tensor-Train barely at all,
since at large $n$ their parameter count is dominated by the router rather
than by the experts.

\begin{table}[t]
\centering
\caption{Parameter counts for $P=113$, expert dimension $m=8$, four router
heads. The dense MLP baseline has $34\,113$ parameters.}
\label{tab:params}
\begin{tabular}{rrrrrrrr}
\toprule
$n$ & $N$ & Unfactorised & MONET-HD & MONET-VD & CP & Tucker & TT \\
\midrule
 4 &   16 &      48\,944 &  14\,948 &  14\,948 & 48\,032 &    30\,880 & 64\,368 \\
 8 &   64 &     188\,544 &  29\,896 &  29\,896 & 52\,672 &    83\,744 & 69\,072 \\
16 &  256 &     739\,712 &  59\,792 &  59\,792 & 61\,952 &   287\,968 & 78\,480 \\
32 & 1024 &  2\,929\,920 & 119\,584 & 119\,584 & 80\,512 & 1\,090\,400 & 97\,296 \\
64 & 4096 & 11\,661\,824 & 239\,168 & 239\,168 & 117\,632 & 1\,121\,376 & 134\,928 \\
\midrule
\multicolumn{2}{r}{growth $\times$} & 238.3 & 16.0 & 16.0 & 2.4 & 36.3 & 2.1 \\
\bottomrule
\end{tabular}
\end{table}

\subsection{Modular addition and its Fourier solution}

The task is to predict $c = (a+b) \bmod P$ from a one-hot encoding of the pair
$(a,b)$, with $P = 113$ prime. There are $P^2 = 12\,769$ examples in total.

A network that generalises rather than memorises is known to implement a
Fourier-multiplication algorithm \citep{nanda2023progress,bussmann2023modadd}.
For a frequency $\omega_k = 2\pi k / P$ it represents each operand by
$\cos(\omega_k a + \varphi_1)$ and $\cos(\omega_k b + \varphi_2)$, and the
readout for class $c$ follows $\cos(\omega_k c + \varphi_1 + \varphi_2)$. The
identity that makes this work is
\begin{equation}
  \sum_{k} \cos\!\big(\omega_k (a + b - c)\big)
  \quad\text{is maximised exactly when}\quad c \equiv a + b \pmod P,
\end{equation}
so a network need only accumulate such terms over a handful of frequencies.

Two consequences shape our methodology. First, the true features of this task
are Fourier frequencies, and there are only $(P-1)/2 = 56$ of them, so
``which feature does this unit represent'' has a definite answer and
monosemanticity can be scored rather than illustrated. Second, since the
operands are one-hot encoded, any weight matrix acting on the input is,
restricted to the block of columns belonging to $a$, literally a function of
$a$; its discrete Fourier transform is therefore directly interpretable. The
same holds for router keys, which under one-hot inputs reduce to learned
lookup tables indexed by the operand. Section~\ref{sec:method} turns these
observations into metrics.
